\subsection{Scelta dello strumento}
% BOZZA: testo di partenza, da rivedere e riscrivere.
Per rispondere alla domanda di ricerca serviva uno strumento in grado di simulare un'intera rete regionale, con centinaia di treni al giorno, e al tempo stesso di rappresentare con precisione i punti in cui la capacità è limitata: i nodi, le stazioni principali e le tratte a binario unico. 
Doveva inoltre permettere di modificare le regole con cui i treni si contendono i binari, perché sono queste regole a determinare come i ritardi si propagano.\\

Gli strumenti professionali più diffusi, come OpenTrack e Trenissimo, simulano la circolazione su una descrizione dettagliata dei binari. Gli autori di railsim ne indicano però i limiti~\cite{kaddoura2024railsim}: il codice è chiuso e non può essere esaminato né modificato, le possibilità di estensione sono ridotte, non è possibile inserire regole proprie per la gestione della circolazione, e l'approccio esclusivamente microscopico richiede dati molto dettagliati su binari e itinerari, il che rende onerosa la simulazione di una rete regionale o nazionale. \\
Esiste anche un simulatore a codice aperto, Flatland, anch'esso microscopico, che usa però una rappresentazione semplificata della realtà ed è impiegato soprattutto per studiare strategie di controllo dei treni con tecniche di apprendimento automatico~\cite{kaddoura2024railsim}. 
All'estremo opposto si trova MATSim usato senza estensioni, che simula il trasporto pubblico a orario ma in cui i veicoli seguono l'orario senza interagire fra loro, e quindi non rappresenta l'esercizio ferroviario~\cite{kaddoura2024railsim}.\\

L'idea iniziale del progetto era lo sviluppo di un simulatore da zero, ma l'idea è stata abbandonata perchè il settore ferroviario comprende molti aspetti legati fra loro, dal moto dei treni all'occupazione dei binari, dalle stazioni agli incroci, e una implementazione scritta da zero non avrebbe raggiunto, nei tempi disponibili, un risultato realistico.\\
% BOZZA: frase su Mesa, da rivedere e riscrivere.
Fra le alternative, c'era anche Mesa~\cite{terhoeven2025mesa}, una libreria Python a codice aperto per la modellazione ad agenti. Mesa fornisce gli elementi generali di un modello ad agenti, ma nessun componente ferroviario: rete, moto dei treni, occupazione dei binari e gestione dell'orario sarebbero stati da scrivere per intero, con le stesse difficoltà dello sviluppo da zero, ma con i vincoli di utilizzare gli elementi della libreria.\\

In seguito ad un analisi di implementazioni in Railsim e la lettura della documentazione, è stato deciso di utilizzare Railsim. \\
Essendo scritto in Java ed open source, può essere esteso con implementazioni specifiche di alcuni comportamenti. Consente di descrivere la stessa rete a livello mesoscopico, microscopico o misto, e quindi di simulare una rete regionale riservando il dettaglio ai punti in cui serve. \\
Infine è sviluppato dalle Ferrovie Federali Svizzere (SBB CHF), il che ne garantisce la serietà del progetto e le sue potenzialità.\\

Lo strumento ha anche due limiti, noti fin dall'inizio. È recente e poco documentato, tanto che per conoscerne il comportamento è stato necessario leggerne il sorgente. Inoltre non costruisce la rete e non stabilisce il livello di dettaglio, che restano a carico di chi lo usa.\\
Nel progetto la rete è stata generata prima a livello mesoscopico, e poi i nodi e i punti più delicati sono stati in seguito descritti a livello microscopico, come illustrato nelle sezioni successive.

\subsection{MATSim}
% BOZZA: testo di partenza, da rivedere e riscrivere.
MATSim (\textit{Multi-Agent Transport Simulation}) è un ambiente di simulazione ad agenti per i sistemi di trasporto, scritto in Java e distribuito con licenza libera~\cite{horni2016matsim}. Il progetto nasce alla fine degli anni Novanta dal lavoro di Kai Nagel, allora al Politecnico federale di Zurigo, a cui si è unito Kay W. Axhausen. È pensato per scenari di grande scala, con milioni di agenti, e per questo ogni suo modello è ridotto all'essenziale.

In MATSim ogni persona della popolazione studiata è un agente, e ogni agente possiede un \textit{piano}: la sequenza delle attività di una giornata e degli spostamenti che le collegano, ognuno con orario, mezzo e percorso. L'unità di analisi è la giornata singola. La domanda di trasporto non è quindi assegnata come flusso fra zone, ma emerge dalla somma dei piani dei singoli agenti.

Una esecuzione è composta da un numero configurabile di iterazioni, ognuna divisa in tre fasi~\cite{horni2016matsim}. Nella prima, la simulazione della mobilità (\textit{mobsim}), tutti i piani scelti vengono eseguiti insieme sulla rete e gli agenti competono per lo spazio disponibile. Nella seconda, la valutazione (\textit{scoring}), ogni piano eseguito riceve un punteggio, che cresce con il tempo dedicato alle attività e diminuisce con il tempo di viaggio e con i ritardi; la funzione usata di norma è quella di Charypar e Nagel~\cite{charypar2005generating}. Nella terza, la ripianificazione (\textit{replanning}), una parte degli agenti, spesso il 10\%, copia il piano scelto e ne modifica un aspetto, come l'orario di partenza, il percorso o il mezzo, mentre gli altri scelgono fra i piani già in memoria in base al punteggio. Ogni agente conserva un numero limitato di piani, cinque nella configurazione predefinita, e scarta quello con il punteggio peggiore. Il ciclo si ripete finché il punteggio medio della popolazione si stabilizza. Si tratta di un algoritmo co-evolutivo: ogni agente migliora i propri piani mentre tutti gli altri fanno lo stesso, e il sistema tende a un equilibrio in cui nessuno può migliorare cambiando da solo il proprio piano.

Il motore di simulazione predefinito, chiamato QSim, usa un modello a code~\cite{gawron1998iterative} e avanza a passi di tempo fissi, di un secondo nella configurazione predefinita. La rete è un grafo orientato di nodi e collegamenti (\textit{link}). Ogni link è una coda in cui il primo veicolo entrato è il primo a uscire, ed è descritto da tre grandezze: il tempo di percorrenza a velocità libera, la capacità di deflusso, cioè quanti veicoli possono uscire nell'unità di tempo, e la capacità di accumulo, cioè quanti veicoli il link può contenere. Un veicolo lascia un link solo quando è trascorso il tempo di percorrenza a velocità libera, la capacità di deflusso del passo corrente non è esaurita e il link successivo ha spazio. Il modello non rappresenta la posizione del veicolo lungo il link, né l'accelerazione, la frenata o la distanza dal veicolo che precede: è il prezzo pagato per poter simulare scenari molto grandi.

Ogni azione che avviene nella simulazione produce un \textit{evento} con il suo istante: un agente parte, un veicolo entra in un link o ne esce, un'attività inizia. Gli eventi sono l'unico risultato grezzo della simulazione, e ogni misura, dai tempi di viaggio ai ritardi, si ottiene elaborandoli. I dati di ingresso di uno scenario minimo sono tre file in formato XML: la configurazione, la rete e la popolazione con i piani iniziali.

Il supporto al trasporto pubblico è stato aggiunto in un secondo momento~\cite{rieser2010adding} e richiede due file ulteriori. Il primo è l'orario (\textit{transit schedule}), che contiene le fermate, ognuna associata a un solo link della rete, e le linee. Ogni linea ha uno o più percorsi, e ogni percorso elenca i link attraversati, le fermate servite, con i tempi di arrivo e di partenza espressi come scarto dalla partenza iniziale, e le singole corse, ognuna con il proprio orario e il veicolo che la effettua. Il secondo file descrive i veicoli e i loro tipi, con capienza, lunghezza e velocità massima. Lo stesso veicolo può essere assegnato a più corse. Durante la simulazione ogni veicolo è guidato da un agente conducente e percorre i link della rete come gli altri veicoli, fermandosi alle fermate per far salire e scendere i passeggeri.

MATSim è costruito come un insieme di moduli sostituibili, collegati fra loro con un meccanismo di iniezione delle dipendenze: chi lo estende può sostituire un componente con una propria implementazione senza modificare il codice esistente. Le estensioni mantenute insieme al nucleo sono chiamate \textit{contrib}, e railsim è una di queste.

Il modello a code è nato per il traffico stradale. Applicato ai treni, non rappresenta ciò che determina la capacità di una ferrovia: l'uso esclusivo di un binario, la distanza di sicurezza fra due treni, i tempi di accelerazione e di frenata, l'incrocio sul binario unico. Per questo railsim sostituisce, per i soli treni, il modo in cui i veicoli si muovono sulla rete, lasciando invariato tutto il resto di MATSim~\cite{kaddoura2024railsim}.

\subsection{Railsim}
Come anticipato nell'introduzione, per l'implementazione del modello, il lavoro si è basato sull'utilizzo di \textit{railsim}~\cite{kaddoura2024railsim}, un'estensione ferroviaria della piattaforma di simulazione ad agenti MATSim.\\
% BOZZA da qui in poi: testo di partenza, da rivedere e riscrivere.

\subsubsection{Obiettivi}
Gli scopi dichiarati di railsim sono tre: verificare se un orario è eseguibile su una data infrastruttura, individuare conflitti e colli di bottiglia, e fare da base per lo studio delle perturbazioni e delle strategie di gestione della circolazione~\cite{kaddoura2024railsim}. 

\subsubsection{Integrazione con MATSim}
railsim non modifica il nucleo di MATSim: aggiunge al motore QSim un motore proprio, che prende in carico i veicoli di un modo di trasporto prefissato, per impostazione predefinita \texttt{rail}. Questi veicoli sono esclusi dal modello a code e si muovono secondo le regole di railsim; tutto il resto resta invariato~\cite{kaddoura2024railsim}. \\
I treni sono quindi veicoli del trasporto pubblico di MATSim: le corse, le fermate, i tempi e il veicolo assegnato a ogni corsa sono letti dall'orario, e ogni treno è guidato da un agente conducente, che railsim sostituisce con una propria versione. 
Lo stesso veicolo può effettuare più corse nella giornata e in questo caso resta sulla rete fra una corsa e la successiva.\\

Tre caratteristiche distinguono il moto di un treno da quello di un veicolo stradale. La prima è l'estensione nello spazio: un treno può occupare più link contemporaneamente, con la testa su uno e la coda su un altro, e un link torna disponibile solo dopo il passaggio della coda. La seconda è la dinamica: il tempo di percorrenza di un link non è fisso, ma dipende dall'accelerazione e dalla frenata del treno. La terza è la prenotazione: un treno può avanzare solo sui tratti di infrastruttura che gli sono stati assegnati, e due treni non possono ottenere lo stesso tratto.

\subsubsection{Infrastruttura: link e risorse}
La rete resta quella di MATSim, un grafo di nodi e link a senso unico, ma gli attributi stradali di capacità e numero di corsie non vengono usati. La velocità massima di un link è l'attributo ordinario di velocità libera; gli altri dati sono attributi propri di railsim, riassunti nella \cref{tab:railsim-link}.

\begin{table}[H]
	\centering
	\caption{Attributi dei link usati da railsim.}
	\label{tab:railsim-link}
	\begin{tabularx}{\textwidth}{lX}
		\toprule
		Attributo & Significato \\
		\midrule
		\texttt{railsimTrainCapacity} & Numero di treni che possono trovarsi insieme sul link. Se manca vale 1. \\
		\texttt{railsimResourceId} & Identificativo della risorsa a cui il link appartiene. I link con lo stesso identificativo si prenotano insieme. \\
		\texttt{railsimMinimumTime} & Tempo minimo fra due treni successivi sul deviatoio alla fine del link. Se manca vale 0. \\
		\texttt{railsimEntry} & Link di ingresso di una stazione: qui il treno può chiedere un percorso alternativo. \\
		\texttt{railsimExit} & Link di uscita di una stazione: è la destinazione del percorso alternativo. \\
		\bottomrule
	\end{tabularx}
\end{table}

Gli attributi si scrivono dentro l'elemento del link, nel file della rete. Il \cref{lst:railsim-capacita} mostra un link che può contenere tre treni insieme.

\begin{lstlisting}[language=XML,caption={Link con capienza di tre treni.},label={lst:railsim-capacita}]
<link id="A" from="A" to="A0" length="500" freespeed="2.77"
      capacity="3600.0" permlanes="1" oneway="1" modes="rail">
  <attributes>
    <attribute name="railsimTrainCapacity"
               class="java.lang.Integer">3</attribute>
  </attributes>
</link>
\end{lstlisting}

Il concetto centrale è quello di \textit{risorsa}. La capacità non è una proprietà del singolo link ma della risorsa, che è formata da uno o più link e ha una capienza pari alla minima fra quelle dei suoi link; un link senza identificativo è una risorsa a sé. Con le risorse si descrivono le situazioni tipiche della ferrovia. Un binario percorribile nei due sensi è una coppia di link di verso opposto con la stessa risorsa: un treno che ne occupa uno blocca anche l'altro. Una sezione di blocco fisso è una sequenza di link con la stessa risorsa e capienza 1. Il blocco mobile si ottiene con link corti, ognuno con la propria risorsa. Un'intersezione fra due binari è protetta dal nodo comune, che resta bloccato finché un treno lo impegna.

Il \cref{lst:railsim-binario-unico} mostra il caso del binario unico: due link di verso opposto fra gli stessi nodi, con lo stesso identificativo di risorsa.

\begin{lstlisting}[language=XML,caption={Binario unico percorribile nei due sensi.},label={lst:railsim-binario-unico}]
<link id="A_B" from="A" to="B" length="200.0" freespeed="50"
      capacity="3600.0" permlanes="1" oneway="1" modes="rail">
  <attributes>
    <attribute name="railsimResourceId"
               class="java.lang.String">AB</attribute>
  </attributes>
</link>
<link id="B_A" from="B" to="A" length="200.0" freespeed="50"
      capacity="3600.0" permlanes="1" oneway="1" modes="rail">
  <attributes>
    <attribute name="railsimResourceId"
               class="java.lang.String">AB</attribute>
  </attributes>
</link>
\end{lstlisting}

Senza l'attributo, i due link sono risorse distinte e descrivono un doppio binario, uno per senso di marcia.

L'infrastruttura può essere descritta a due livelli di dettaglio, che differiscono per ciò che un link rappresenta. \\
railsim è in grado di gestirli entrambi utilizzando lo stesso motore.

Nella descrizione microscopica ogni link è un singolo binario e ogni risorsa ha capienza 1: su un binario può trovarsi un solo treno alla volta. È il livello in cui si distinguono le sezioni di blocco, i singoli binari di una stazione e i collegamenti fra di essi. Una stazione microscopica è formata da un link per ogni binario di banchina, ciascuno con la propria fermata. Tutte le fermate appartengono alla stessa area di fermata, e in questo modo railsim sa che si tratta di binari alternativi della stessa stazione; i link che portano alla stazione e quelli che ne escono sono marcati come ingresso e uscita.

Nella descrizione mesoscopica un link rappresenta un fascio di binari paralleli, e la sua capienza indica quanti treni possono percorrerlo insieme. La disposizione dei binari non è descritta: ne resta solo il numero. Un link mesoscopico può iniziare o finire soltanto in un punto in cui un treno può effettivamente cambiare binario. Una stazione mesoscopica è una risorsa con capienza pari al numero dei suoi binari: railsim sa quanti treni può contenere, ma non su quale binario si trova ciascuno.

I due livelli hanno costi diversi. Quello microscopico richiede la disposizione esatta dei binari e più tempo di calcolo; a quello mesoscopico bastano le lunghezze e il numero dei binari. Possono essere usati insieme nella stessa rete, descrivendo in dettaglio l'area di interesse e in modo mesoscopico il resto~\cite{kaddoura2024railsim}.

In entrambi i casi railsim si limita a simulare la rete che riceve: non la costruisce e non fornisce i dati dei binari. Nel progetto la rete mesoscopica è stata generata a partire dall'orario e dalle misure dei tracciati, mentre la descrizione microscopica delle stazioni principali, binario per binario, è stata costruita appositamente, come illustrato nelle sezioni successive.

Il \cref{lst:railsim-meso} mostra una tratta mesoscopica che ammette due treni per senso di marcia, e il \cref{lst:railsim-uscita} un link di uscita da una stazione.

\begin{lstlisting}[language=XML,caption={Tratta mesoscopica con capienza di due treni per senso.},label={lst:railsim-meso}]
<link id="l_bc_F" from="n_b2" to="n_c1" length="4000"
      freespeed="27.777" capacity="3600.0" permlanes="1" modes="rail">
  <attributes>
    <attribute name="railsimTrainCapacity"
               class="java.lang.Integer">2</attribute>
  </attributes>
</link>
<link id="l_bc_R" from="n_c1" to="n_b2" length="4000"
      freespeed="27.777" capacity="3600.0" permlanes="1" modes="rail">
  <attributes>
    <attribute name="railsimTrainCapacity"
               class="java.lang.Integer">2</attribute>
  </attributes>
</link>
\end{lstlisting}

\begin{lstlisting}[language=XML,caption={Link di uscita da una stazione.},label={lst:railsim-uscita}]
<link id="hl" from="H" to="L" length="3000" freespeed="25"
      capacity="3600.0" permlanes="1" modes="rail">
  <attributes>
    <attribute name="railsimExit"
               class="java.lang.Boolean">true</attribute>
  </attributes>
</link>
\end{lstlisting}

\subsubsection{Treni e dinamica del moto}
Ogni tipo di treno è descritto dalla lunghezza, dalla velocità massima, dalla capienza e da due attributi propri di railsim: l'accelerazione e la decelerazione, in metri al secondo quadrato. Se mancano, valgono i valori predefiniti della configurazione. 
Il \cref{lst:railsim-treno} mostra la definizione di un tipo di treno nel file dei veicoli: il modo di trasporto \texttt{rail} è quello che lo affida al motore di railsim.

\begin{lstlisting}[language=XML,caption={Definizione di un tipo di treno.},label={lst:railsim-treno}]
<vehicleType id="fvDosto">
  <attributes>
    <attribute name="railsimAcceleration"
               class="java.lang.Double">0.4</attribute>
    <attribute name="railsimDeceleration"
               class="java.lang.Double">0.5</attribute>
  </attributes>
  <capacity seats="606" standingRoomInPersons="0"/>
  <length meter="200.0"/>
  <width meter="1.435"/>
  <maximumVelocity meterPerSecond="55.556"/>
  <networkMode networkMode="rail"/>
</vehicleType>
\end{lstlisting}

La velocità di un treno risulta da quattro elementi: le sue prestazioni, la sua velocità massima, la velocità massima dei link e la disponibilità dell'infrastruttura davanti a lui. \\
\begin{itemize}
    \item Un treno accelera finché i tratti successivi sono liberi e non ha raggiunto la velocità consentita.
    \item Frena quando una risorsa che ha richiesto è occupata da un altro treno, oppure quando su un link successivo la velocità consentita è inferiore alla sua.
    \item La velocità da mantenere è deciso dal \textit{profilo di velocità}. 
    \begin{itemize}
        \item Il profilo predefinito porta il treno alla massima velocità
        \item Un secondo profilo regola la velocità in base alla posizione del treno e all'orario del prossimo arrivo.
    \end{itemize}
\end{itemize}

\subsubsection{Prenotazione del percorso}
Un treno non occupa solo i link corrente ma anche l'infrastruttura davanti a sé. \\
La lunghezza del tratto prenotato è calcolata a ogni aggiornamento ed è pari allo spazio che serve al treno per fermarsi alla velocità corrente. \\
Con velocità $v$ e decelerazione $d$ vale
\begin{equation}
	s = \frac{v^2}{2d}
\end{equation}
Il tratto prenotato non va mai oltre la fermata successiva finché il treno non riparte, la prenotazione si interrompe alla stazione. \\
Se una risorsa richiesta non è disponibile, il treno riceve un permesso più corto, rallenta fino a fermarsi prima del tratto occupato e ripete la richiesta dopo un intervallo fisso. In questo modo la distanza di sicurezza fra due treni non è un parametro, ma nasce dalla velocità e dalla capacità di frenata di chi segue.

\subsubsection{Gestione della circolazione}
Le decisioni su chi può prenotare che cosa sono affidate a due componenti, definiti come interfacce e sostituibili da chi estende lo strumento. \\
Il primo, \texttt{TrainDisposition}, riceve le richieste di prenotazione dei treni e risponde con la distanza concessa e la velocità consentita. \\
La versione predefinita concede le risorse nell'ordine in cui vengono chieste, senza priorità fra i treni. Si occupa anche delle deviazioni: quando un treno raggiunge un link di ingresso e uno dei link che dovrebbe percorrere fino al primo link di uscita è occupato e senza capienza residua, viene cercato un percorso alternativo fra l'ingresso e l'uscita. È il meccanismo che permette a un treno di fermarsi a un binario diverso da quello previsto nella stessa stazione.\\
Il secondo componente, \texttt{DeadlockAvoidance}, serve a evitare gli stalli, cioè le situazioni in cui due o più treni aspettano ciascuno una risorsa tenuta dall'altro. La versione predefinita, secondo la sua stessa documentazione, evita molti casi ma non garantisce l'assenza di stalli. \\
La documentazione di railsim distingue un caso locale, due treni di verso opposto su una tratta con capienza condivisa, e un caso globale, in cui lo stallo nasce da una sequenza di tratte con capienze diverse, e descrive il secondo come un problema ancora aperto. Gli stessi autori indicano le strategie di gestione della circolazione come la parte più difficile e ancora in sviluppo~\cite{kaddoura2024railsim}.

\subsubsection{Eventi}
Per restare compatibile con MATSim, railsim produce gli eventi ordinari, come l'ingresso di un veicolo in un link, e ne aggiunge quattro propri, elencati nella \cref{tab:railsim-eventi}.

\begin{table}[H]
	\centering
	\caption{Eventi aggiunti da railsim.}
	\label{tab:railsim-eventi}
	\begin{tabularx}{\textwidth}{lX}
		\toprule
		Evento & Contenuto \\
		\midrule
		\texttt{RailsimTrainLeavesLinkEvent} & La coda del treno ha lasciato un link. \\
		\texttt{RailsimLinkStateChangeEvent} & Un link è passato da libero a occupato o viceversa, con il treno che lo occupa. \\
		\texttt{RailsimTrainStateEvent} & Stato del treno: posizione della testa e della coda, velocità, velocità obiettivo, accelerazione. È prodotto a ogni cambiamento di stato e a intervalli regolari. \\
		\texttt{RailsimDetourEvent} & Il treno è stato deviato, con la parte di percorso modificata. \\
		\bottomrule
	\end{tabularx}
\end{table}

\subsubsection{Configurazione}
I parametri generali di railsim sono riportati nella \cref{tab:railsim-config}, con i valori predefiniti.

\begin{table}[H]
	\centering
	\caption{Parametri di configurazione di railsim.}
	\label{tab:railsim-config}
	\begin{tabularx}{\textwidth}{llX}
		\toprule
		Parametro & Predefinito & Significato \\
		\midrule
		\texttt{networkModes} & \texttt{rail} & Modi di trasporto gestiti dal motore di railsim. \\
		\texttt{accelerationDefault} & \SI{0,5}{\metre\per\second\squared} & Accelerazione usata se il tipo di treno non la indica. \\
		\texttt{decelerationDefault} & \SI{0,5}{\metre\per\second\squared} & Decelerazione usata se il tipo di treno non la indica. \\
		\texttt{pollInterval} & \SI{10}{\second} & Attesa prima di ripetere una richiesta di prenotazione respinta. \\
		\texttt{updateInterval} & \SI{10}{\second} & Intervallo massimo fra due eventi di stato di un treno. \\
		\texttt{speedProfile} & \texttt{maxSpeed} & Profilo di velocità: \texttt{maxSpeed}, \texttt{adaptive} oppure \texttt{custom}. \\
		\texttt{vehicleCirculation} & \texttt{yes} & Circolazione dello stesso veicolo su più corse. \\
		\bottomrule
	\end{tabularx}
\end{table}

\subsubsection{Che cosa resta a carico di chi lo usa}
railsim simula il moto e i conflitti dei treni su una rete e con un orario che riceve già pronti. Non costruisce la rete e non stabilisce quale treno effettua ogni corsa né su quale binario ferma, perché legge queste informazioni dall'orario. \\
Non prevede depositi o binari di ricovero: un treno fermo fra due corse occupa il binario su cui si trova. \\
Non calcola indicatori di puntualità né il consumo di energia, e non introduce guasti o perturbazioni casuali, per cui lo stesso scenario produce sempre lo stesso risultato. \\
Come si è visto, infine, la prevenzione degli stalli è solo parziale. Le sezioni che seguono descrivono come il progetto ha affrontato ciascuno di questi punti.

\subsection{Confronto con Mesa e Mesa-Geo}
Come anticipato, un'alternativa a Railsim sarebbe potuta essere Mesa che si presenta come la controparte in Python degli ambienti classici per la modellazione ad agenti, come NetLogo, Repast e MASON.\\
Mette a disposizione le classi di base per il modello e per gli agenti, gli spazi in cui collocarli, come griglie e reti, la raccolta dei dati e una visualizzazione nel browser~\cite{terhoeven2025mesa}.
È uno strumento generale, che non fa ipotesi sul fenomeno simulato: il comportamento degli agenti e le regole con cui interagiscono sono interamente a carico di chi scrive il modello.
La sua estensione Mesa-Geo~\cite{wang2022mesageo} aggiunge i dati geografici. Gli agenti hanno una geometria, cioè un punto, una linea o un poligono, con il proprio sistema di riferimento, e possono essere creati direttamente da file geografici. Per questo progetto sarebbe stata utile nella prima fase, quella in cui stazioni e tracciati vengono ricavati dai dati cartografici e mostrati su una mappa. 
La geografia però non è la parte difficile del problema: lo sono le regole dell'esercizio ferroviario, che né Mesa né Mesa-Geo contengono.

La \cref{tab:confronto-strumenti} riassume il confronto. Con Mesa, anche esteso con Mesa-Geo, la parte ferroviaria sarebbe stata da scrivere per intero; con railsim è già disponibile, e il lavoro si è potuto concentrare sulla costruzione della rete e sugli aspetti che lo strumento non prevede.

\begin{table}[H]
	\centering
	\caption{Confronto fra Mesa, Mesa-Geo e railsim rispetto alle esigenze del progetto.}
	\label{tab:confronto-strumenti}
	\begin{tabularx}{\textwidth}{lXXX}
		\toprule
		Aspetto & Mesa & Mesa-Geo & railsim \\
		\midrule
		Natura & libreria generale per modelli ad agenti & estensione geografica di Mesa & estensione ferroviaria di MATSim \\
		Linguaggio & Python & Python & Java \\
		Spazio & griglie e reti & geometrie con sistema di riferimento & grafo di link con risorse e capienze \\
		Dati geografici & non previsti & lettura da file geografici & coordinate dei nodi \\
		Moto dei treni & da scrivere & da scrivere & accelerazione, frenata, lunghezza del treno \\
		Occupazione dei binari & da scrivere & da scrivere & prenotazione delle risorse \\
		Orario & da scrivere & da scrivere & orario del trasporto pubblico di MATSim \\
		\bottomrule
	\end{tabularx}
\end{table}